% Estrutura baseada em Wohlin et al. (escopo e planejamento do experimento).
\chapter[Planejamento do Experimento]{Planejamento do Experimento}
\label{cap:planejamento}

\section{Definição do escopo do experimento}
\label{sec:escopo}

\subsection{Objetivo}
% {aqui vai o objetivo do estudo no template GQM, porque define o que será medido}
% (analisar / com o propósito de / com respeito à / do ponto de vista de / no contexto de)

\subsection{Questão de pesquisa}
% {aqui vão as questões do estudo empírico, porque guiam a coleta e a análise}

\subsection{Trabalhos relacionados}
% {aqui vão os estudos da revisão mais próximos do nosso e a diferença deles para o nosso, porque posiciona a contribuição}

\section{Planejamento}
\label{sec:planejamento}

\subsection{Contexto}
% {aqui vão os projetos open source web ágeis escolhidos, os critérios de seleção e as histórias de usuário usadas, porque o contexto do GQM exige esse cenário}
% Objetos experimentais, população e amostra (histórias do CAPJu).

\subsection{Variáveis do experimento}
% {aqui vão as variáveis independentes (ex.: LLM, técnica de prompting), dependentes (ex.: acurácia dos critérios) e as hipóteses nula e alternativa, porque definem o que o experimento testa}

\subsubsection*{Fator experimental}
% Estratégia de contexto, com dois níveis (sem RAG e com RAG).

\subsubsection*{Variável independente}
% TODO

\subsubsection*{Variáveis dependentes}
% Acurácia (Acc@Hit e Acc@Cor); similaridade como medida complementar.

\subsubsection*{Variáveis controladas}
% LLM e versão, temperatura, prompt, corpus, top-k.

\subsection{Formulação das hipóteses}
% H0 e H1, nível de significância.

\subsection{Desenho do experimento}
% {aqui vão os tratamentos e a forma de alocação, porque determina como os dados serão comparados}
% Um fator, dois tratamentos, comparação pareada.

\subsection{Teste de normalidade}
% Shapiro-Wilk sobre as diferenças pareadas.

\subsection{Teste de hipóteses}
% {aqui vão as métricas e os métodos de análise (estatística descritiva, testes), porque a escolha deve vir antes da coleta}
% Teste t pareado ou Wilcoxon.

\subsection{Instrumentação}
% {aqui vão as ferramentas, os prompts, o gabarito de referência e os scripts, porque permite reproduzir o estudo}

\subsection{Procedimento de aplicação dos instrumentos}
% {aqui vão os passos de execução em ordem, porque a etapa do TCC2 vai seguir esse roteiro}

\subsection{Validade do experimento}
% {aqui vão as ameaças por tipo e as mitigações, porque é parte obrigatória do planejamento nos TCCs anteriores}

\subsubsection*{Validade de construto}
% {aqui vai se as medidas representam de fato a qualidade dos critérios, porque é o principal risco em métricas de texto}

\subsubsection*{Validade interna}
% {aqui vão fatores que podem distorcer causa e efeito (ex.: não determinismo dos LLMs), porque afeta a confiança nos resultados}

\subsubsection*{Validade externa}
% {aqui vai até onde os resultados generalizam para outros projetos e modelos, porque o contexto é restrito a apps web open source}

\subsubsection*{Validade de conclusão}
% TODO

\subsection{Análise complementar}
% Medidas complementares (ex.: similaridade de cosseno), de caráter exploratório.
